\subsection{Promoting AI Development Policies that Discourage Authoritarian Use}
\label{sec:6.4.3}
The technical safeguards above raise the cost of building a system that can be turned toward surveillance or control. Whether anyone has to pay that cost is a policy question, and the part of it that belongs here is narrow: what rules govern which AI capabilities get built, sold, and deployed, regardless of who is coordinating with whom.

Three levers close three different doors. Export and trade controls decide who can buy, and the EU's 2021 recast of its dual-use export regulation added a human-rights-based control for cyber-surveillance items, requiring an exporter who becomes aware that an item is likely to be used in serious human-rights violations to seek approval even where the item falls outside the pre-existing control lists \autocite{europeanunion2021regulation}. That rule is narrower than it sounds, since neither “internal repression” nor “serious violation” is defined in the text and member states have read both inconsistently; what it establishes is that a license rather than a sale is where this kind of misuse can be stopped. Prohibited-practices regulation decides what can be built at all, by naming specific designs off-limits regardless of who is building them, and section~\ref{sec:10.4} takes up the instance. Developer-facing standards decide what a builder treats as normal, which is the weakest lever and the fastest to arrive: the Toronto Declaration set out non-discrimination standards for machine-learning systems years before most governments had draft legislation on the subject \autocite{amnestyinternational2018toronto}.

None of the three works alone. Export controls with no prohibited-practices list at home let a government ban a technology abroad while permitting it domestically, and a design standard with no legal force behind it binds only the developers already inclined to follow it. What makes them work together is that an authoritarian government, or a company willing to serve one, has to get through all three rather than through whichever one happens to apply in its jurisdiction.
